\section{Total Ward extraction through field degree four}
\label{tkw-section}

The following argument completes the low-degree \emph{total-kernel}
classification. It does not impose gauge invariance on individual
activity coefficients. Its exact statement concerns subsequential
infinite-lattice kernels; finite-period winding terms and localization
of the canceled functional remain separate questions.

\subsection{Hypotheses and the reference action}
\label{tkw-hypotheses}

For a compact connected simple matrix group, use Hermitian Lie-algebra
coordinates and a fixed invariant positive trace form. Consider analytic
total corrections \(E_N\) on exhausting compatible cubic tori. Assume
the local gauge identity, output-unit translations, full signed
permutations and reality on real fields established for the actual
local sector in Section~\ref{msy-section}. Assume connected activity
representations on a fixed cell partition with
\begin{equation}
 \sup_I\sum_{H\ni I}e^{a_0|H|}\|E_{N,H}\|_b\le\epsilon .
 \label{tkw-activity}
\end{equation}
Here \(a_0,b>0\), the cell size \(J\), and the scalar coordinate count
\(D\) per cell are fixed uniformly in volume and preceding refinement.
The total symmetries need not hold for the chosen coefficients.

Let \(S_N\) be a real analytic, exactly gauge-invariant and spatially
invariant reference, with \(S_N(0)=DS_N(0)=0\), and with uniform
fixed-order kernel moment bounds through degree four. Assume its
limiting midpoint Hessian has the normalization
\begin{equation}
 \Gamma_2^S(p)=|p|^2I-pp^T+O(|p|^4).
 \label{tkw-reference}
\end{equation}
The normalized Wilson action is one choice. The actual homogeneous
classical action \(S_{MR}\) is another: its Hessian is \(q_{MR}\),
whose uniform normalization was proved by the alias formula in
Notebook07, and the classical strict-support construction supplies
the higher fixed-order bounds. Ellipticity alone would not establish
\eqref{tkw-reference}.

\begin{theorem}[Full marginal completion from the total Ward identity]
\label{tkw-theorem}
Under these assumptions, choose a joint subsequential infinite-lattice
limit of the kernels of \(E_N,S_N\) through degree four. There is a real
coefficient \(\beta\) such that the midpoint Fourier vertices of the
residual \(F=E-\beta S\) satisfy
\begin{align}
 \Gamma_2^F(p)&=O(|p|^4),\notag\\
 \Gamma_3^F(p_1,p_2,p_3)
     &=O\!\left(\max_j|p_j|^3\right),\notag\\
 \Gamma_4^F(p_1,p_2,p_3,p_4)
     &=O\!\left(\max_j|p_j|^2\right),
 \qquad \sum_jp_j=0.
 \label{tkw-remainders}
\end{align}
Thus the total Taylor coefficients of field degree \(n=2,3,4\)
and momentum degree at most \(4-n\) agree with \(\beta S\).
The one-point kernel vanishes; the vacuum is unrestricted.
\end{theorem}

For a semisimple product there is one coefficient per simple ideal.
The analogous statement subtracts gauge-invariant references with the
individual normalized Maxwell blocks. A single-coupling rebase is
immediate for a simple group, or if those coefficients coincide.
Reality can be omitted if a complex coefficient \(\beta\) is permitted.

\subsection{Fixed-order kernels and their joint limit}
\label{tkw-kernel-limit}

Write \(K_n=D^nE_N(0)\). Coordinate Cauchy on one coefficient gives
\begin{equation}
 \sum_{i_1,\ldots,i_n\in H}
  |D^nE_{N,H}(0)[e_{i_1},\ldots,e_{i_n}]|
 \le (D|H|+n)^n b^{-n}\|E_{N,H}\|_b .
 \label{tkw-coordinate-Cauchy}
\end{equation}
Indeed the exact monomial count multiplied by \(n!\) is
\(n!\binom{D|H|+n-1}{n}\), bounded by the displayed factor.
Fixing the first coordinate only decreases the sum. The graph
Steiner length of slots contained in connected \(H\) is at most
\(|H|-1\). For \(0<\sigma<a_0\), therefore,
\begin{equation}
 \sup_{i_1}\sum_{i_2,\ldots,i_n}
 e^{\sigma\tau(i_1,\ldots,i_n)}
 |K_n(i_1,\ldots,i_n)|
 \le b^{-n}\epsilon
       \sup_{m\ge1}(Dm+n)^n e^{-(a_0-\sigma)m}.
 \label{tkw-kernel-bound}
\end{equation}
Scalar and color indices are included in the sums. A further fixed
exponent loss absorbs every fixed spatial moment. Conversion from
the \(J\)-cell metric to the unit-site metric retains \(J\)-dependent
constants.

Global constant gauges imply \(K_1=0\) separately at each bond,
because its covector annihilates every Lie commutator. Translation
invariance permits an anchor at the origin. The bounds then give a
diagonal subsequence of all fixed displacements, simultaneously for
\(n\le4\) and for the reference. Fatou preserves exponential
summability. Permutation symmetry of the slots, signed permutations,
translation and global adjoint covariance pass to this limit.
The differentiated Ward identities below pass as well: their
coefficient equations contain finite nearest-neighbor shifts and
lower kernels with coinciding slots. Equivalently one may first
test against finitely supported fields and gauges.

The limiting Fourier vertices are analytic near the simultaneous
origin, and their fixed derivatives are justified by the spatial
moments. No unnormalized infinite-volume total functional is being
summed; even a nonzero vacuum density would preclude that.

\subsection{The exact logarithmic gauge generator}
\label{tkw-generator-section}

For \(U_b=e^{iA_b}\), use
\[
 U_b(t)=e^{it\lambda_{b_-}}U_be^{-it\lambda_{b_+}},
 \qquad (d\lambda)_b=\lambda_{b_+}-\lambda_{b_-}.
\]
Differentiating the exponential chart, with
\(\operatorname{ad}_A X=[A,X]\), gives
\begin{equation}
 \delta_\lambda A
 =i[\lambda_-,A]-f(i\operatorname{ad}_A)d\lambda,
 \qquad
 f(z)=\frac{z}{1-e^{-z}}
       =1+\frac z2+\frac{z^2}{12}+O(z^4).
 \label{tkw-generator}
\end{equation}
In particular
\begin{align}
 \delta_\lambda A
   &=-d\lambda+L_\lambda(A)+Q_\lambda(A,A)
                        +O(A^4\lambda),\notag\\
 L_\lambda(A)_b
   &=\frac i2[\lambda_{b_-}+\lambda_{b_+},A_b],\notag\\
 Q_\lambda(A,B)_b
   &=\frac1{24}\bigl(
      [A_b,[B_b,d\lambda_b]]+
      [B_b,[A_b,d\lambda_b]]\bigr).
 \label{tkw-LQ}
\end{align}
The cubic Bernoulli coefficient is zero. Comparing powers in
\(DE(A)[\delta_\lambda A]=0\), using \(K_1=0\), gives
\begin{align}
 K_2(A,d\lambda)&=0,\notag\\
 K_3(A,A,d\lambda)&=2K_2(A,L_\lambda(A)),\notag\\
 K_4(A,A,A,d\lambda)
   &=3K_3(A,A,L_\lambda(A))
       +6K_2(A,Q_\lambda(A,A)).
 \label{tkw-Ward-diagonal}
\end{align}
Polarization yields
\begin{align}
 K_3(A,B,d\lambda)
   &=K_2(A,L_\lambda(B))+K_2(B,L_\lambda(A)),\label{tkw-Ward-cubic}\\
 K_4(A,B,C,d\lambda)
   &=K_3(A,B,L_\lambda(C))
      +K_3(A,C,L_\lambda(B))
      +K_3(B,C,L_\lambda(A))\notag\\
   &\quad+2K_2(A,Q_\lambda(B,C))
      +2K_2(B,Q_\lambda(A,C))
      +2K_2(C,Q_\lambda(A,B)).
 \label{tkw-Ward-quartic}
\end{align}
These are total-kernel identities. Subtracting just a quadratic
polynomial would generally violate them. Subtracting a multiple
of an exactly invariant reference preserves them.

The generator and differentiated identities are independently
derived here. Their published counterparts are Ba{\l}aban,
CMP 109 (1987), (4.9)--(4.15), pp.~283--284; the full-action
subtraction is (4.38)--(4.45), pp.~291--292. The stronger
coefficientwise orbit hypotheses of that article are not assumed.

\subsection{Midpoint phases and parity}
\label{tkw-midpoint-section}

Use bond-midpoint Fourier fields and site Fourier gauges:
\[
 A_\mu(x)=\int e^{ip\cdot(x+e_\mu/2)}a_\mu(p)\,dp,
 \qquad \lambda(x)=\int e^{ik\cdot x}\lambda(k)\,dk .
\]
On \(\sum_jp_j=0\), anchoring the first bond gives the vertex phase
\begin{equation}
 \exp\!\left\{i\sum_{j=2}^n
   p_j\cdot(x_j+e_{\mu_j}/2-e_{\mu_1}/2)\right\}.
 \label{tkw-vertex-phase}
\end{equation}
Only local analyticity near zero momentum is required; midpoint
vertices need not be individually periodic under every \(2\pi\)
shift. The Fourier coefficients in \eqref{tkw-LQ} are exactly
\begin{align}
 (d\lambda)_\mu&:\quad
       2i\sin(k_\mu/2)\lambda(k),\notag\\
 L_\lambda(A)_\mu&:\quad
       i\cos(k_\mu/2)[\lambda(k),a_\mu(p)],\notag\\
 Q_\lambda(A,B)_\mu&:\quad
       \frac{i}{12}\sin(k_\mu/2)
       \bigl([a_\mu(p),[b_\mu(q),\lambda(k)]]
           +[b_\mu(q),[a_\mu(p),\lambda(k)]]\bigr).
 \label{tkw-generator-phases}
\end{align}
Their output momenta are \(k,p+k,p+q+k\). Thus \(L\) has
a bounded analytic coefficient, while \(Q\) supplies one
power of gauge momentum.

Full inversion sends \(a_\mu(p)\) to \(-a_\mu(-p)\).
Consequently
\begin{equation}
 \Gamma_n(-p_1,\ldots,-p_n)
        =(-1)^n\Gamma_n(p_1,\ldots,p_n).
 \label{tkw-parity}
\end{equation}
The cubic vertex is odd; the quadratic and quartic vertices
are even. Midpoint phases are essential for this form.

\subsection{Proof of the marginal completion}
\label{tkw-completion-proof}

\begin{proof}[Proof of Theorem~\ref{tkw-theorem}]
Theorem~\ref{msy-Maxwell-theorem} gives
\[
 \Gamma_2^E(p)=\beta(|p|^2I-pp^T)+O(|p|^4).
\]
For completeness, the continuous limiting Ward relation
\(\Gamma_2(p)\widehat p=0\) first implies \(\Gamma_2(0)=0\)
by taking \(p=te_\mu\). Reflections and permutations restrict
the degree-two tensor to diagonal
\(A p_\mu^2+B\sum_{\nu\ne\mu}p_\nu^2\) and off-diagonal
\(C p_\mu p_\nu\). Transversality gives \(A=0,C=-B\).
Put \(\beta=B\), which is real under the stated reality
assumption. It can be extracted as
\begin{equation}
 \beta=\frac12\,
  \partial_{p_\nu}^2(\Gamma_2^E)_{\mu\mu}(0),
       \qquad \mu\ne\nu,
 \label{tkw-beta}
\end{equation}
in a trace-orthonormal color direction. The reference
normalization proves \(\Gamma_2^F(p)=O(|p|^4)\).

Apply \eqref{tkw-Ward-cubic} to the residual and momenta
\(tp,tq,tk\), with \(p+q+k=0\). Its right side is \(O(t^4)\):
each residual Hessian is evaluated at small momentum, and
the \(L\) coefficient in \eqref{tkw-generator-phases} is bounded.
Parity gives
\[
 \Gamma_3^F(tp,tq,tk)=tP_1(p,q,k)+O(t^3).
\]
The left side contracts its last vector index with
\(2i\sin(tk/2)=itk+O(t^3)\). Its degree-two coefficient
therefore vanishes. Slot permutation symmetry gives, for
each choice of colors,
\begin{equation}
 p_\mu(P_1)_{\mu\nu\rho}=0,\qquad
 q_\nu(P_1)_{\mu\nu\rho}=0,\qquad
 k_\rho(P_1)_{\mu\nu\rho}=0.
 \label{tkw-linear-transverse}
\end{equation}

Eliminate \(k=-p-q\) and write
\[
 (P_1)_{\mu\nu\rho}
  =\sum_\sigma A_{\mu\nu\rho;\sigma}p_\sigma
                   +B_{\mu\nu\rho;\sigma}q_\sigma .
\]
In the first relation of \eqref{tkw-linear-transverse},
the independent mixed monomials \(p_\mu q_\sigma\)
force every \(B_{\mu\nu\rho;\sigma}\) to vanish.
The second relation then forces every \(A_{\mu\nu\rho;\sigma}\)
to vanish by the independent monomials \(q_\nu p_\sigma\).
Thus \(P_1=0\), proving the cubic remainder. This elementary
argument permits arbitrary complex/color coefficients and
does not presume a separately invariant cubic functional.

Finally apply \eqref{tkw-Ward-quartic}, scaling all four
momenta by \(t\). Its \(K_3,L\) terms are \(O(t^3)\);
its \(K_2,Q\) terms are \(O(t^5)\), because \(Q\) has
the sine factor. By parity write
\[
 \Gamma_4^F(tp_1,\ldots,tp_4)=P_0+O(t^2).
\]
The order-\(t\) term on the left contracts \(P_0\) with
the last momentum \(p_4\). That momentum is arbitrary:
the other three can always sum to \(-p_4\). Hence
every component of \(P_0\) vanishes. This proves the
quartic remainder and completes the theorem.
\end{proof}

The argument is an identity of the integrated vertices.
It does not choose a unique density modulo lattice divergences.

\subsection{Reference independence and the full action}
\label{tkw-full-action-section}

For the standard Wilson reference, the marginal continuum
polynomial is explicit. With Hermitian coordinates set
\[
 f_{\mu\nu}=\partial_\mu A_\nu-\partial_\nu A_\mu,
 \qquad F_{\mu\nu}=f_{\mu\nu}+i[A_\mu,A_\nu].
\]
Its engineering-degree-four part is
\begin{equation}
 \frac12\sum_{\mu<\nu}\tau(F_{\mu\nu}^2)
 =\frac12\sum_{\mu<\nu}\tau(f_{\mu\nu}^2)
   +i\sum_{\mu<\nu}\tau(f_{\mu\nu}[A_\mu,A_\nu])
   -\frac12\sum_{\mu<\nu}\tau([A_\mu,A_\nu]^2).
 \label{tkw-Wilson-polynomial}
\end{equation}
Indeed bond angles \(tA\) evaluated at midpoints of a
lattice of spacing \(t\) give the plaquette
\(P=I+it^2F_{\mu\nu}+O(t^3)\). Therefore
\(1-\frac12\tau(P+P^{-1})
 =\frac12t^4\tau(F_{\mu\nu}^2)+O(t^5)\).
This verifies the normalization and the relative cubic
and quartic coefficients.

No explicit Wilson expansion is needed in the theorem.
Applying it to the difference of two exactly invariant
normalized references shows that their marginal total
vertices agree. In particular one may subtract the actual
homogeneous classical reference directly:
\begin{equation}
 e^{-g^{-2}S_{MR}+E}
  =e^{-(g^{-2}-\beta)S_{MR}+(E-\beta S_{MR})}.
 \label{tkw-coupling-rebase}
\end{equation}
The vacuum is separate. This identity neither proves
positivity of \(g^{-2}-\beta\) nor the required sign,
size or accumulated drift of its actual increments.

\subsection{Moment remainders and the remaining norm question}
\label{tkw-scope-section}

For the residual anchored kernels use the \emph{bond midpoints}
\(\xi_j=x_j+e_{\mu_j}/2\), and define
\begin{equation}
 T_{n,m}=\sup_{i_1}\sum_{i_2,\ldots,i_n}
  \left(\sum_{j=2}^n|\xi_j-\xi_1|_1\right)^m
                  |K_n^F(i_1,\ldots,i_n)|.
 \label{tkw-moments}
\end{equation}
They are bounded by the moments of \(E\) plus
\(|\beta|\) times those of \(S\). The phase
\eqref{tkw-vertex-phase} and the real one-variable
Taylor integral remainder give componentwise
\begin{align}
 |\Gamma_2^F(p)|
   &\le \frac{T_{2,4}}{4!}\|p\|_\infty^4,\notag\\
 |\Gamma_3^F(p_1,p_2,p_3)|
   &\le \frac{T_{3,3}}{3!}
                  \max_{j\ge2}\|p_j\|_\infty^3,\notag\\
 |\Gamma_4^F(p_1,p_2,p_3,p_4)|
   &\le \frac{T_{4,2}}{2!}
                  \max_{j\ge2}\|p_j\|_\infty^2 .
 \label{tkw-quantitative}
\end{align}
Fixed vector/color row sums add only their fixed dimensions.
The exponential kernel bound makes these moments finite.

The cancellations also hold in ordinary tail-anchored
coordinates: conversion multiplies each midpoint component
by a nonzero analytic phase and cannot lower its vanishing
order. Thus all residual ordinary momentum jets through
degrees \(3,2,1\) vanish for field degrees \(2,3,4\),
respectively. This includes the weaker anchored moments
used in second-order quadratic, first-order cubic and
zeroth-order quartic spatial subtractions.

The constants can depend strongly on the fixed old grouping
\(J\). Since the common-domain construction chooses \(J\)
after \(R\), an estimate \(C(J)R^{-\alpha}\) is not
automatically a contraction. No such claim is made.
The remaining tasks are localization of the canceled
total without false coefficientwise Ward assumptions,
finite-period wrapping control or a consistent infinite
correction, a spatial-gain norm with the correct parameter
order, and the actual relevant/marginal recursion.
The theorem proves the total-kernel Taylor classification;
it does not supply these separate conclusions.
